\documentclass[10pt,a4paper,twocolumn]{article}

\usepackage[a4paper,top=1.75cm,bottom=2.0cm,left=1.75cm,right=1.75cm,columnsep=0.65cm]{geometry}
\usepackage{times}
\usepackage[T1]{fontenc}
\usepackage[utf8]{inputenc}
\usepackage{microtype}
\usepackage{amsmath}
\usepackage{booktabs}
\usepackage{tabularx}
\usepackage{enumitem}
\usepackage{xcolor}
\usepackage{natbib}
\usepackage[hidelinks]{hyperref}
\usepackage{titlesec}

\setlength{\parindent}{1em}
\setlength{\parskip}{0pt}
\setlength{\textfloatsep}{8pt plus 2pt minus 2pt}
\setlength{\floatsep}{7pt plus 2pt minus 2pt}
\setlength{\intextsep}{7pt plus 2pt minus 2pt}
\setlist{nosep,leftmargin=1.4em}
\titleformat{\section}{\large\bfseries}{\thesection}{0.55em}{}
\titleformat{\subsection}{\normalsize\bfseries}{\thesubsection}{0.45em}{}
\titlespacing*{\section}{0pt}{1.4ex plus .4ex}{0.7ex}
\titlespacing*{\subsection}{0pt}{1.0ex plus .3ex}{0.4ex}

\title{Reproducing LLMLingua: Evaluating Accuracy--Efficiency Trade-offs in Long-Context Prompt Compression}

\author{
CHEN IAT HONG \quad MC654102\\
KONG KIN MENG \quad MC652817\\
CHE CHON HONG \quad MC551841\\
\textit{University of Macau}\\
\textit{CISC7021 Applied Natural Language Processing -- Course Project Final Report}
}

\date{}

\begin{document}
\maketitle
\vspace{-0.8cm}

% ===========================================================================
% SKELETON ONLY: headings are fixed; every section, subsection and
% subsubsection body is replaced by the placeholder line
%     (text text text)
% Replace each placeholder and its "% TODO" note with real prose, and keep the
% heading structure unchanged.
% ===========================================================================

\begin{abstract}
(text text text)

% TODO: 150--200 words. Motivation (long prompts cost time and money, position
% bias), what was reproduced (uncompressed, head-and-tail truncation, LLMLingua,
% LLMLingua-enhanced), the setup (LongBench multifieldqa_en, 150 samples,
% Llama-3.2-3B-Instruct, 2x/4x), the headline numbers, and the one-sentence
% conclusion (compression trades quality for latency; LLMLingua did not beat
% truncation at an equal token budget in this setting).
\end{abstract}

% ---------------------------------------------------------------------------
\section{Introduction}
(text text text)

\subsection{Motivation}
(text text text)

% TODO: why long-context prompting is expensive; redundancy and lost-in-the-middle;
% why prompt compression (no retraining, model-agnostic) is attractive; why a
% resource-constrained, reproducible evaluation is still missing.

\subsection{Literature Review}
(text text text)

% TODO: condense report/acl_latex.tex Sections 2--4 into 1--2 paragraphs:
% prompt-compression survey, Selective Context, LLMLingua, LongLLMLingua,
% lost-in-the-middle, LongBench. Cite with \citet / \citep.

\subsection{Research Questions}
(text text text)

% TODO: restate RQ1 (ratio vs accuracy), RQ2 (LLMLingua vs truncation at the same
% token budget), RQ3 (end-to-end latency including compression overhead).

\subsection{Contributions}
(text text text)

% TODO: bullet list — (i) shared harness with a fixed sample index, prompt
% template and metrics; (ii) four compression variants re-implemented and run on
% the full 150-sample split; (iii) seven-setting comparison with paired
% significance tests; (iv) a single-machine latency rerun; (v) a negative result
% for LLMLingua under an equal-budget comparison.

% ---------------------------------------------------------------------------
\section{Detailed Description}
(text text text)

\subsection{Task and Dataset}
(text text text)

% TODO: LongBench multifieldqa_en; question + retrieved documents + gold answers;
% why all 150 official test examples were kept; how the fixed index is stored in
% shared/samples.json (seed 42).

\subsection{Models and Hardware}
(text text text)

% TODO: target model Llama-3.2-3B-Instruct (gated id vs unsloth mirror); 128k
% native context and why a 32k prompt budget was self-imposed; compressor models
% (microsoft/phi-2; the enhanced variant in Part 4); the four machines and the
% single-machine rerun used for latency.

\subsection{Shared Evaluation Harness}
(text text text)

% TODO: one contract for all members — files, what is frozen, what a member may
% change. Expand on the four sub-blocks below.

\subsubsection{Sample Index}
(text text text)

% TODO: 150 ids, order, seed, why cross-part comparability depends on it.

\subsubsection{Prompt Template}
(text text text)

% TODO: the chat template and instruction/context formatting; why prompt_tokens
% counts the chat-templated prompt; the token-budget fallback.

\subsubsection{Metrics}
(text text text)

% TODO: SQuAD-style Exact Match and token-level F1; mean over 150 samples; what
% EM/F1 ignore (e.g. LaTeX fragments, truncation at 64 new tokens).

\subsubsection{Configuration}
(text text text)

% TODO: greedy decoding, max_new_tokens = 64, MAX_PROMPT_TOKENS, dtype, seeds.

\subsection{Compared Methods}
(text text text)

% TODO: one paragraph per method plus the target ratio definitions (document
% ratio vs whole-prompt ratio).

\subsubsection{Part 1: Uncompressed Baseline}
(text text text)

% TODO: mechanism, parameters, expected role as reference.

\subsubsection{Part 2: Head-and-Tail Truncation}
(text text text)

% TODO: budget b = round(n / r); head_frac = 0.5; instruction and question are
% never truncated; why this is a non-semantic baseline.

\subsubsection{Part 3: LLMLingua}
(text text text)

% TODO: coarse-to-fine, iterative token filtering, rates 0.5 / 0.25, phi-2 as
% compressor, measurement tokenizer.

\subsubsection{Part 4: LLMLingua-Enhanced}
(text text text)

% TODO: what changed relative to Part 3 (sentence-level filtering, question-aware
% selection, 4-bit quantisation), and why the improvement cannot be attributed
% to a single factor.

\subsection{Implementation}
(text text text)

% TODO: how the pipeline is structured end to end; the three sub-blocks below.

\subsubsection{Two-Stage Compress and Evaluate Pipeline}
(text text text)

% TODO: compressor and target model never co-resident; two environments because
% of the transformers version conflict; JSONL checkpoints and resume.

\subsubsection{Chunked Compression and Chunked Prefill}
(text text text)

% TODO: 1,400-token compressor chunks in Part 3; --chunk-size / --chunk-threshold
% prefill in Part 2; how chunking is recorded per row.

\subsubsection{Memory and Out-of-Memory Handling}
(text text text)

% TODO: Colab T4 limits, OOM retry with halved prompt budget, skipped-and-retried
% rows, oom_skipped_ids.

\subsubsection{System Overview}
(text text text)

% TODO: insert the pipeline / component diagram here (figure with caption and
% label). One figure covering data -> (compression) -> generation -> scoring.

\subsection{Experimental Setup}
(text text text)

% TODO: which settings were run, on which machine, at which ratio, when, and
% which result file each row comes from.

\subsection{Evaluation Protocol}
(text text text)

% TODO: primary metrics; paired-by-id F1 differences; 95% CI = mean +/- 1.96 SE;
% how "significant" is decided; which contrasts are reported; why per-part
% latencies are not compared across machines.

% ---------------------------------------------------------------------------
\section{Results and Evaluation}
(text text text)

\subsection{Main Results}
(text text text)

% TODO: the seven-setting table (uncompressed, head-tail 2x/4x, LLMLingua
% 2x/4x, LLMLingua-enhanced 2x/4x): n, EM, F1, prompt tokens, token reduction,
% compression / generation / total latency. Reference part4 results_table.md.

\subsection{RQ1: Accuracy versus Compression Ratio}
(text text text)

% TODO: monotone F1 decline from 1x to 2x to 4x in all three compressed families;
% include the accuracy--compression figure (figures/accuracy_compression_tradeoff).

\subsection{RQ2: LLMLingua versus Truncation at Equal Token Budget}
(text text text)

% TODO: paired contrasts — LLMLingua loses to truncation (-0.1440 / -0.1414);
% enhanced LLMLingua improves over LLMLingua (+0.0444 / +0.0587) but still loses
% to truncation (-0.0995 / -0.0827); all CIs exclude zero.

\subsection{RQ3: End-to-End Latency}
(text text text)

% TODO: use the single-machine rerun only; truncation cuts total latency by
% ~44% / ~66%; LLMLingua 2x is slower than uncompressed once compression is
% charged; Part 4 has no common-hardware rerun (state this explicitly).

\subsection{Qualitative Analysis}
(text text text)

% TODO: reproducible success and failure cases from
% part3_llmlingua/qualitative_cases.md; what compression removes that the
% question needed; what survives intact.

\subsection{Discussion}
(text text text)

% TODO: why perplexity-based token removal underperforms a dumb baseline here
% (fragmentation, names/numbers, 3B target model, short prompts, question not
% used for budget allocation); what this implies about the paper's claims.

\subsection{Threats to Validity}
(text text text)

% TODO: per-part hardware differences; stale results_same_machine export;
% differing verdict for Uncompressed -> Head-tail 2x between runs; single
% dataset, single model, single head_frac; five-sample pilots are not final.

% ---------------------------------------------------------------------------
\section{Conclusion}
(text text text)

\subsection{Summary of Findings}
(text text text)

% TODO: answer RQ1--RQ3 in three sentences, each with its number.

\subsection{Lessons Learned}
(text text text)

% TODO: about reproduction, about evaluation contracts, about measuring
% end-to-end cost rather than token counts.

\subsection{Problems and Difficulties}
(text text text)

% TODO: gated model access; Windows console encoding crash at sample 87; the
% compressor context limit forcing chunking; Colab T4 OOM forcing chunked
% prefill; the Windows commit-limit failure that blocked Part 4 locally.

\subsection{Limitations}
(text text text)

% TODO: scope and measurement limits, one bullet per limitation.

\subsection{Future Work}
(text text text)

% TODO: answer-position failure analysis; LongLLMLingua-style reordering;
% larger target models; more ratios and datasets; latency measurement on one
% common GPU for all four parts.

% ---------------------------------------------------------------------------
\section{Use of Generative AI}
(text text text)

% TODO: which tools were used for what (literature organisation, drafting,
% formatting, code assistance), and how numbers and claims were verified.

% ---------------------------------------------------------------------------
% References: carried over from report/acl_latex.tex. Keep only the entries that
% are actually cited, and add any new ones in the same format.
\begin{thebibliography}{9}

\bibitem[Bai et al.(2024)]{bai-etal-2024-longbench}
Yushi Bai, Xin Lv, Jiajie Zhang, Hongchang Lyu, Jiankai Tang, Zhidian Huang, Zhengxiao Du, Xiao Liu, Aohan Zeng, Lei Hou, Yuxiao Dong, Jie Tang, and Juanzi Li. 2024.
\newblock LongBench: A bilingual, multitask benchmark for long context understanding.
\newblock In \emph{Proceedings of the 62nd Annual Meeting of the Association for Computational Linguistics (Volume 1: Long Papers)}, pages 3119--3137.

\bibitem[Jiang et al.(2023)]{jiang-etal-2023-llmlingua}
Huiqiang Jiang, Qianhui Wu, Chin-Yew Lin, Yuqing Yang, and Lili Qiu. 2023.
\newblock LLMLingua: Compressing prompts for accelerated inference of large language models.
\newblock In \emph{Proceedings of the 2023 Conference on Empirical Methods in Natural Language Processing}, pages 13358--13376.

\bibitem[Jiang et al.(2024)]{jiang-etal-2024-longllmlingua}
Huiqiang Jiang, Qianhui Wu, Xufang Luo, Dongsheng Li, Chin-Yew Lin, Yuqing Yang, and Lili Qiu. 2024.
\newblock LongLLMLingua: Accelerating and enhancing LLMs in long context scenarios via prompt compression.
\newblock In \emph{Proceedings of the 62nd Annual Meeting of the Association for Computational Linguistics (Volume 1: Long Papers)}, pages 1658--1677.

\bibitem[Kwiatkowski et al.(2019)]{kwiatkowski-etal-2019-natural}
Tom Kwiatkowski, Jennimaria Palomaki, Olivia Redfield, Michael Collins, Ankur Parikh, Chris Alberti, Danielle Epstein, Illia Polosukhin, Jacob Devlin, Kenton Lee, Kristina Toutanova, Llion Jones, Matthew Kelcey, Ming-Wei Chang, Andrew M. Dai, Jakob Uszkoreit, Quoc Le, and Slav Petrov. 2019.
\newblock Natural Questions: A benchmark for question answering research.
\newblock \emph{Transactions of the Association for Computational Linguistics}, 7:453--466.

\bibitem[Li et al.(2023)]{li-etal-2023-compressing}
Yucheng Li, Bo Dong, Frank Guerin, and Chenghua Lin. 2023.
\newblock Compressing context to enhance inference efficiency of large language models.
\newblock In \emph{Proceedings of the 2023 Conference on Empirical Methods in Natural Language Processing}, pages 6342--6353.

\bibitem[Li et al.(2025)]{li-etal-2025-prompt-survey}
Zongqian Li, Yinhong Liu, Yixuan Su, and Nigel Collier. 2025.
\newblock Prompt compression for large language models: A survey.
\newblock In \emph{Proceedings of the 2025 Conference of the Nations of the Americas Chapter of the Association for Computational Linguistics: Human Language Technologies (Volume 1: Long Papers)}, pages 5807--5823.

\bibitem[Liu et al.(2024)]{liu-etal-2024-lost}
Nelson F. Liu, Kevin Lin, John Hewitt, Ashwin Paranjape, Michele Bevilacqua, Fabio Petroni, and Percy Liang. 2024.
\newblock Lost in the middle: How language models use long contexts.
\newblock \emph{Transactions of the Association for Computational Linguistics}, 12:157--173.

\end{thebibliography}

\end{document}
